\documentclass[11pt,a4paper]{article}
\usepackage[utf8]{inputenc}
\usepackage[T1]{fontenc}
\usepackage{geometry}
\geometry{margin=1in}
\usepackage{hyperref}
\usepackage{enumitem}
\usepackage{booktabs}
\usepackage{longtable}
\usepackage{xcolor}
\usepackage{titlesec}
\usepackage{amsmath,amssymb}

\titleformat{\section}{\Large\bfseries\color{blue!70!black}}{}{0em}{}
\titleformat{\subsection}{\large\bfseries\color{blue!50!black}}{}{0em}{}
\titleformat{\subsubsection}{\normalsize\bfseries\color{blue!30!black}}{}{0em}{}

\title{\textbf{Replication Plan}\\[0.5em]
\large Electronic Specific Heat Capacities and Entropies from Density Matrix Quantum Monte Carlo Using Gaussian Process Regression\\[0.3em]
\normalsize OSTI ID: 1993311}
\author{Auto-generated Replication Blueprint}
\date{\today}

\begin{document}
\maketitle
\tableofcontents
\newpage

%---------------------------------------------------------------------
\section{Paper Overview}
%---------------------------------------------------------------------
This paper uses density matrix quantum Monte Carlo (DMQMC) and its initiator variant (PIP-DMQMC) to compute finite-temperature properties (internal energy, specific heat capacity, entropy) of small molecular and model systems (e.g., uniform electron gas, small atoms/molecules). Because QMC energies at finite temperature are noisy, the authors apply Gaussian process (GP) regression to smooth the energy-vs-temperature data and then analytically differentiate the GP to obtain the specific heat $C_v(T) = dE/dT$ and integrate to obtain the entropy. They benchmark against exact full configuration interaction (FCI) results.

%---------------------------------------------------------------------
\section{Detailed Workflow}
%---------------------------------------------------------------------

\subsection{Phase 1: Exact Diagonalization Benchmarks}
\begin{enumerate}[label=\textbf{1.\arabic*}]
  \item \textbf{Implement or use FCI} for small systems (few electrons in small basis sets).
  \item \textbf{Compute exact $E(\beta)$} (energy vs.\ inverse temperature) via the thermal density matrix:
    \begin{equation*}
      E(\beta) = \frac{\text{Tr}[H e^{-\beta H}]}{\text{Tr}[e^{-\beta H}]}
    \end{equation*}
  \item \textbf{Compute exact $C_v(\beta)$ and $S(\beta)$} from the partition function.
\end{enumerate}

\subsection{Phase 2: DMQMC / PIP-DMQMC Simulations}
\begin{enumerate}[label=\textbf{2.\arabic*}]
  \item \textbf{Set up the Hamiltonian} in second quantization (molecular integrals or UEG parameters).
  \item \textbf{Run DMQMC.}
    \begin{itemize}
      \item Stochastic propagation of the density matrix in imaginary time.
      \item Sample at multiple inverse temperatures $\beta$ values.
      \item Record $E(\beta)$ with statistical error bars.
    \end{itemize}
  \item \textbf{Run PIP-DMQMC} (interaction-picture variant) for improved efficiency.
\end{enumerate}

\subsection{Phase 3: Gaussian Process Regression}
\begin{enumerate}[label=\textbf{3.\arabic*}]
  \item \textbf{Fit GP to $E(\beta)$ data.}
    \begin{itemize}
      \item Input: $\{(\beta_i, E_i, \sigma_i)\}$ from DMQMC.
      \item Kernel: squared-exponential (RBF) or Mat\'{e}rn.
      \item Heteroscedastic noise model using the QMC error bars.
    \end{itemize}
  \item \textbf{Optimize GP hyperparameters} (length scale, signal variance) via marginal likelihood.
  \item \textbf{Compute derivatives analytically.}
    \begin{itemize}
      \item $C_v(\beta) = -\beta^2 \frac{dE}{d\beta}$; the GP derivative is analytic.
      \item Entropy via thermodynamic integration.
    \end{itemize}
  \item \textbf{Propagate uncertainties} through the GP to get error bands on $C_v$ and $S$.
\end{enumerate}

\subsection{Phase 4: Validation}
\begin{enumerate}[label=\textbf{4.\arabic*}]
  \item Compare GP-derived $C_v(\beta)$ and $S(\beta)$ against exact FCI results.
  \item Reproduce convergence of DMQMC results with walker number.
  \item Reproduce all figures: $E(\beta)$, $C_v(\beta)$, $S(\beta)$ with error bands.
\end{enumerate}

%---------------------------------------------------------------------
\section{Datasets}
%---------------------------------------------------------------------

\begin{longtable}{p{4.5cm} p{3.5cm} p{6cm}}
\toprule
\textbf{Dataset / Source} & \textbf{Type} & \textbf{Location / Notes} \\
\midrule
\endfirsthead
\toprule
\textbf{Dataset / Source} & \textbf{Type} & \textbf{Location / Notes} \\
\midrule
\endhead
Molecular integrals & Generated & From PySCF or similar quantum chemistry code \\
\midrule
UEG parameters & Specified & Electron number, $r_s$, basis set size from paper \\
\midrule
DMQMC energy data & Generated & Output of QMC simulations \\
\bottomrule
\end{longtable}

%---------------------------------------------------------------------
\section{Models, Tools, and Software}
%---------------------------------------------------------------------

\begin{longtable}{p{4.5cm} p{2.5cm} p{6.5cm}}
\toprule
\textbf{Tool / Library} & \textbf{Version} & \textbf{Purpose / URL} \\
\midrule
\endfirsthead
\toprule
\textbf{Tool / Library} & \textbf{Version} & \textbf{Purpose / URL} \\
\midrule
\endhead
HANDE & latest & QMC code (DMQMC, FCIQMC) \newline \url{https://github.com/hande-qmc/hande} \\
\midrule
PySCF & $\geq 2.0$ & Molecular integrals, FCI benchmarks \newline \url{https://pyscf.org/} \\
\midrule
scikit-learn & $\geq 1.0$ & Gaussian process regression \newline \url{https://scikit-learn.org/} \\
\midrule
GPy & latest & Alternative GP library \newline \url{https://sheffieldml.github.io/GPy/} \\
\midrule
NumPy / SciPy & $\geq 1.21$ & Numerical operations \\
\midrule
Matplotlib & $\geq 3.4$ & Plotting \\
\midrule
pyblock & latest & QMC error analysis (blocking/reblocking) \newline \url{https://github.com/jsspencer/pyblock} \\
\bottomrule
\end{longtable}

\subsection{Hardware Requirements}
\begin{itemize}
  \item DMQMC: HPC cluster for large walker populations; small systems feasible on workstation.
  \item GP regression: any laptop.
\end{itemize}

\end{document}
